\section{Experimental Setup}
\label{sec:protocol}

\subsection{Data and model}

We used 75M tokens of FineWeb-Edu, a collection of educational web pages
\cite{penedo2024fineweb}, tokenized with the GPT-2 tokenizer. The documents were
split into training, validation, and test sets so that no document appears in
more than one set, and the split was kept fixed throughout the study; the test
set was not used. Every run reads the training tokens in the same order, so a
200M-token run passes over the training set about 2.7 times. The model is a
small GPT-2-style Transformer with 8 layers, a width of 384, six attention
heads, and a context of 512 tokens (Table~\ref{tab:experimental-protocol}).

\subsection{Training}

Each training step processed 32,768 tokens (64 sequences of 512 tokens). The
learning rate rose linearly over the first 5\% of training and then fell along a
cosine curve to 10\% of its peak. We clipped the gradient norm at 1.0 and
trained in mixed precision. Each run in the main experiment used one NVIDIA A10
GPU.

\begin{table*}[t]
\caption{Setup of the main experiment.}
\label{tab:experimental-protocol}
\centering
\small
\begin{tabular}{p{0.12\textwidth} p{0.22\textwidth} p{0.56\textwidth}}
\hline
Category & Setting & Value \\
\hline
Model & Architecture & GPT-2-style decoder-only Transformer: 8 layers, width 384, 6 attention heads, MLP width 1,536, pre-layer normalization, GELU, learned positions, tied input and output embeddings \\
Model & Tokenizer and context & GPT-2 tokenizer (50,257 tokens); 512-token context \\
Data & Corpus & FineWeb-Edu; fixed snapshot recorded in the repository \\
Data & Split & 75M training, 5M validation, and 5M test tokens; no document in more than one set \\
Training & Length & 200M tokens (6,103 steps; 199,983,104 tokens processed) \\
Training & Batch & 32,768 tokens per step (2 $\times$ 32 sequences of 512 tokens) \\
Training & Learning-rate schedule & 5\% linear warm-up, then cosine decay to 10\% of the peak; gradient clipping at 1.0 \\
Training & Precision & Mixed precision (bfloat16, with TF32 matrix multiplication); Newton--Schulz iteration in float32 \\
Evaluation & Loss curve & Every $\approx$10M tokens on a fixed 999,424-token part of the validation set \\
Evaluation & Final loss & Full validation set (4,997,120 tokens) \\
Hardware & GPU & One NVIDIA A10 per run \\
Seeds & Main experiment & 301, 302, 303, 304, 305 \\
\hline
\end{tabular}
\end{table*}


\subsection{Choosing the learning rate}
\label{sec:selection}

The two optimizers share every setting except SOAP--Muon's extra step
(Table~\ref{tab:optimizer-definitions}), so the learning rate was the only
setting we tuned. Tuning used shorter, 40M-token runs with a separate seed
(101), each scored by the area under its validation-loss curve (defined in
Section~\ref{sec:measurements}).

The first round of tuning (Table~\ref{tab:screening-record}) selected
SOAP--Muon on both MLP matrices at a learning rate of $1.2\times10^{-3}$. In
hindsight, this round was unbalanced. The baseline's best rate was the largest
one we tried, so a larger rate might have been better, and the SOAP--Muon
candidates changed only the learning rate of the MLP matrices, so the two grids
were not comparable. After an earlier five-seed comparison at this setting
(Section~\ref{sec:sensitivity}), we ran a second round in which both optimizers
use the same learning rate for every parameter
(Table~\ref{tab:shared-rate-screen}). Both did best at $4.8\times10^{-3}$ and
worse at half and at double that rate, so the main experiment uses
$4.8\times10^{-3}$ for both. Appendix~\ref{app:exploratory} discusses the
tuning runs further.

% Generated by scripts/generate_exploratory_tables.py; do not edit manually.
\begin{table*}[t]
\caption{First round of learning-rate tuning (seed 101, 40M-token runs). Each entry is the area under the validation-loss curve between 10M and 30M tokens, divided by 40M; lower is better. For Moonlight Muon, the learning rate applies to all parameters; for SOAP--Muon, it applies only to the MLP matrices that use the SOAP step, with all other parameters at $1.2\times10^{-3}$. The Moonlight Muon values above $1.2\times10^{-3}$ were added after the earlier comparison, and Table~\ref{tab:shared-rate-screen} gives the second round. Dashes mark settings that were not run.}
\label{tab:screening-record}
\centering
\scriptsize
\begin{tabular}{l r r r r r r}
\hline
Optimizer & $3\times10^{-4}$ & $6\times10^{-4}$ & $1.2\times10^{-3}$ & $2.4\times10^{-3}$ & $4.8\times10^{-3}$ & $9.6\times10^{-3}$ \\
\hline
Moonlight Muon & 2.953 & 2.776 & 2.636 & 2.549 & 2.518 & 2.544 \\
SOAP--Muon, up only & -- & 2.641 & 2.630 & 2.623 & -- & -- \\
SOAP--Muon, down only & -- & 2.611 & 2.620 & 2.635 & -- & -- \\
SOAP--Muon, both & -- & 2.624 & 2.610 & 2.631 & -- & -- \\
\hline
\end{tabular}
\end{table*}

% Generated by scripts/generate_exploratory_tables.py; do not edit manually.
\begin{table}[t]
\caption{Post-confirmation matched shared-rate screen on seed 101. Every eligible Muon, SOAP, and AdamW group uses the displayed rate; SOAP--Muon preconditions both MLP projections. Entries are token-normalised validation-NLL areas over 10M, 20M, and 30M tokens; lower is better. This single-seed screen selects a configuration but is not confirmation evidence.}
\label{tab:shared-rate-screen}
\centering
\footnotesize
\begin{tabular}{l r r r}
\hline
Method & $2.4\times10^{-3}$ & $4.8\times10^{-3}$ & $9.6\times10^{-3}$ \\
\hline
Moonlight Muon & 2.549 & 2.518 & 2.544 \\
SOAP--Muon, both & 2.537 & 2.507 & 2.526 \\
\hline
\end{tabular}
\end{table}


\subsection{Main experiment}

The main experiment trains each optimizer with five new seeds (301--305) that
were not used for tuning. A seed fixes the model's initial weights, and every
run sees the data in the same order, so the two runs with the same seed differ
only in the optimizer; we call them a pair. The SOAP--Muon run for seed 305 was
first placed on a different GPU model (an A10G), so we repeated it on an A10 to
keep each pair on identical hardware. The rule for this repeat was written down
before its result was known, and Section~\ref{sec:sensitivity} shows how much
the choice matters.

\subsection{Measurements}
\label{sec:measurements}

The validation loss is the average negative log-likelihood of the next token on
held-out text; lower is better. We measured it every 10M training tokens on a
fixed 1M-token part of the validation set and, at the end of training, on the
full 5M-token validation set.

A single final loss shows only where a run ended, not how quickly it got there.
Our main measure, chosen before the main experiment, therefore summarizes the
whole loss curve: we join the measurements from 10M to 200M tokens with
straight lines and divide the area under this curve by the 200M-token training
length. The result is, in effect, the run's average validation loss during
training. (It is scaled by about 0.95 because the curve starts at 10M rather
than 0 tokens; the factor is the same for every run, so it does not affect
comparisons.) A run whose loss falls faster has a smaller area, so the measure
rewards fast learning as well as a low final loss.
Fig.~\ref{fig:nll-difference}(b) shows how the difference between the two
optimizers' areas arises.

We also report the final loss, which an earlier project proposal had named as
the main measure; the time needed to reach a loss of 4.25 and 4.20; the
training speed; and the peak GPU memory. For each measure and seed, we subtract
the Moonlight Muon value from the SOAP--Muon value, so negative differences
favor SOAP--Muon. We report the mean difference over the five seeds, how many
seeds favor SOAP--Muon, and a 95\% confidence interval for the mean based on
the $t$-distribution; the repository also reports bootstrap intervals and sign
tests. With five seeds, these numbers describe the size and consistency of an
effect rather than give a definitive test.
